%--------------------------
\begin{frame}{Question 6}

\begin{block}{}
Une demande internationale ne devrait être déposée que si :

\begin{itemize}
  \item[$\circ$] a. le déposant pense, lors du dépôt, avoir de sérieuses chances d’obtenir des brevets 
  nationaux dans une majorité d’États contractants.   
  \item[$\circ$] b. le déposant est d’avis que la validité de la revendication de priorité n’est pas 
  contestable.   
  \item[$\circ$] c. tous les déposants sont ressortissants d’États contractants du PCT. 
  \item[$\bullet$] d. Aucune des réponses a, b ou c n’est correcte.
\end{itemize}

\end{block}

\end{frame}

\begin{frame}{Question 6}

\begin{itemize}
  \item[$\circ$] \textbf{a.} \; Faux : le PCT est un outil procédural de dépôt international,
  indépendant des chances de délivrance dans la majorité des États.

  \item[$\circ$] \textbf{b.} \; Faux : une revendication de priorité peut être contestée,
  mais cela n’empêche pas le dépôt d’une demande internationale
  (\pctart{8}.1)).

  \item[$\circ$] \textbf{c.} \; Faux : il suffit qu’au moins un déposant ait nationalité ou domicile
  dans un État contractant, pas nécessairement tous
  (\pctrule{18}.3)).

  \item[$\bullet$] \textbf{d.} \; Correct : aucune des propositions a, b ou c n’est exigée pour déposer une demande PCT.
\end{itemize}

\end{frame}